\documentclass[12pt]{article}
\usepackage[T1]{fontenc}
\usepackage[utf8]{inputenc}
\usepackage{lmodern}
\usepackage{textcomp}
\usepackage{enumitem}
\usepackage{geometry}
\geometry{margin=1in}
\begin{document}
\begin{center}
Answer Key - History - K-12 Level Assessment 1 \
\small (Answers are ordered exactly as in the question paper.)
\end{center}

\section*{Multiple Choice}
\begin{enumerate}[label=\arabic*.]
\item \textbf{Answer:} D
\\ \textit{Options:} A) collecting plants; B) fishing; C) agriculture; D) foraging
\\ \textit{Note:} Foraging is the correct answer because it encompasses the primary means of subsistence for early humans, who relied on hunting, gathering, and fishing for sustenance before the advent of agriculture and settled societies.
\item \textbf{Answer:} A
\\ \textit{Options:} A) Mediterranean; B) Egyptian; C) Mesopotamian; D) Indonesian
\\ \textit{Note:} The Minoans inhabited the Mediterranean island of Crete because it was strategically located in the central part of the Mediterranean Sea, which facilitated trade and cultural exchange with other ancient civilizations.
\item \textbf{Answer:} A
\\ \textit{Options:} A) the king's body was mummified and treated as if it were alive; B) hundreds of the elite were sacrificed by ritual decapitation; C) the body was entombed with lavish burial goods within an enormous pyramid; D) all of the above
\\ \textit{Note:} The correct answer highlights the Inca practice of mummifying their kings and treating their preserved bodies with great reverence, as they believed the deceased rulers retained power and influence over their subjects even in death.
\item \textbf{Answer:} A
\\ \textit{Options:} A) sacrificed adults and occasionally children; B) burnt their own houses to the ground; C) fasted for long periods; D) abstained from sexual relations
\\ \textit{Note:} The Minoans engaged in the practice of sacrificing adults and occasionally children as a ritualistic means to appease their gods during times of crisis, reflecting their religious beliefs and cultural practices in ancient society.
\item \textbf{Answer:} B
\\ \textit{Options:} A) They looked and thought like us.; B) They looked like us but did not think like us.; C) They thought like us but did not look like us.; D) They did not look or think like us.
\\ \textit{Note:} The correct answer highlights that while anatomically modern humans shared similar physical characteristics with contemporary humans, they lacked the advanced cognitive abilities and complex thought processes that define modern human behavior and culture.
\end{enumerate}

\section*{True / False}
\begin{enumerate}[label=\arabic*.]
\setcounter{enumi}{5}
\item \textbf{Answer:} True
\\ \textit{Options:} A) True; B) False
\\ \textit{Note:} Many ancient elites prioritized displaying their wealth and resources as a means of showcasing power and status within their societies, often through grand architectural projects, lavish banquets, and luxurious clothing, rather than hoarding their riches.
\item \textbf{Answer:} True
\\ \textit{Options:} A) True; B) False
\\ \textit{Note:} After 400,000 years ago, hominid evolution saw a significant increase in brain size, particularly with the emergence of Homo neanderthalensis and Homo sapiens, which is linked to advancements in tool use, social structures, and cognitive abilities.
\item \textbf{Answer:} False
\\ \textit{Options:} A) True; B) False
\\ \textit{Note:} The statement is false because while agriculture can increase food production, advancements in technology and industrialization have led to fewer people needing to work in farming while still meeting the protein needs of a growing population.
\item \textbf{Answer:} False
\\ \textit{Options:} A) True; B) False
\\ \textit{Note:} The climate of Greater Australia 20,000 years ago included diverse environments, such as temperate forests and glacial regions, due to variations in geography and climate patterns during that time.
\item \textbf{Answer:} True
\\ \textit{Options:} A) True; B) False
\\ \textit{Note:} The construction of canals for irrigation allowed the Caral civilization to effectively manage water resources, which in turn supported agriculture, population growth, and the establishment of a complex social structure.
\end{enumerate}

\section*{Long Form Response}
\begin{enumerate}[label=\arabic*.]
\setcounter{enumi}{10}
\item \textbf{Answer:} Australia plays a significant role in early human history, particularly through the discovery of ancient rock engravings, which are among the oldest forms of artistic expression in the world. These etchings, created by Indigenous Australians, provide valuable insights into the cultural practices, beliefs, and social structures of ancient societies. By studying these artworks, historians and archaeologists can better understand how early humans interacted with their environment and each other. The implications of this evidence extend beyond art; they reveal a rich tapestry of history that highlights the continuity and resilience of Indigenous cultures over thousands of years. Overall, Australia's rock engravings are crucial for comprehending the complexities of early human life and the development of civilization.
\\ \textit{Note:} correct because it highlights the significance of Australia's ancient rock engravings as crucial evidence of early artistic expression, offering insights into the cultural practices and social structures of Indigenous Australians, thereby enhancing our understanding of early human interactions with their environment and communities.
\item \textbf{Answer:} The transition from traditional foraging to maize-based agriculture around 2000 B.P. significantly impacted social structures and daily life in the Southwest. With the adoption of agriculture, communities began to settle in one place, leading to the development of permanent villages. This shift allowed for increased food production, which supported larger populations and led to more complex social hierarchies. As people focused on farming, daily life became centered around agricultural practices, such as planting, harvesting, and managing crops, which influenced family roles and labor division. Overall, this transition marked a pivotal change that shaped the culture and organization of societies in the region.
\\ \textit{Note:} correct because it identifies the transition to maize-based agriculture as a catalyst for permanent settlement, increased food production, population growth, and the emergence of complex social hierarchies, all of which fundamentally reshaped social structures and daily life in the Southwest.
\end{enumerate}

\vfill
\noindent\textit{End of Answer Key}
\end{document}